% Irodov Part V §5.2 87--96, concludes Light Interference.
% Source PDF ~210--211, printed pp.201--202; source answer pp.328--329.
\interferenceproblem{87}{201}
平凸透镜球面半径 $R=40$cm 与玻璃板接触。反射牛顿环中某一条环的初始半径
$r=2.5$mm；缓慢将透镜整体上提
$\Delta h=5.0\,\mathrm{\mu m}$，观察\emph{原来同一阶}环的新半径。
\exsolution
初始膜厚
$h(r)=r^2/(2R)$。
透镜抬高后
$h'(r')=\Delta h+r'^2/(2R)$。
同一干涉阶要求同一膜厚，
$\Delta h+r'^2/(2R)=r^2/(2R)$，
于是
\[
\boxed{r'=\sqrt{r^2-2R\Delta h}}
=\sqrt{(2.5\,\mathrm{mm})^2-
2(0.40\,\mathrm m)(5.0\,\mathrm{\mu m})}
=\boxed{1.5\,\mathrm{mm}}.
\]
\excheck 必须确认根号内非负；再提升到 $2R\Delta h>r^2$ 后原阶环将消失。

\interferenceproblem{88}{201}
一块平凸透镜球面顶部被磨成半径
$r_0=3.0$mm 的平坦区，该平面与下方玻璃板接触。
原球面曲率半径
$R=150$cm。
用波长 $\lambda=655$nm 的光观察反射牛顿环，
求第六个亮环半径。
\exsolution
磨去顶部后，从接触圆边缘 $r_0$ 开始，
膜厚
$h(r)\simeq(r^2-r_0^2)/(2R)$。
反射时只有一个界面产生 $\pi$ 相位反转，
第 $k$ 个亮环满足
$2h=(k-\tfrac12)\lambda$。
因此
\[
\boxed{r_k=\sqrt{r_0^2+(k-\tfrac12)\lambda R}}.
\]
第六亮环
\[
\boxed{r_6=
\sqrt{(3.0\,\mathrm{mm})^2+
5.5(655\,\mathrm{nm})(1.50\,\mathrm m)}
\simeq3.8\,\mathrm{mm}}.
\]
\excheck 不能将平坦接触斑的半径 $r_0$ 当作新的曲率半径；它的平方作为几何偏移进入环半径。

\interferenceproblem{89}{201}
平凸玻璃透镜球面曲率半径
$R=12.5$cm，压在玻璃板上。
反射光第十、第十五暗牛顿环直径分别为
$d_{10}=1.00$mm、
$d_{15}=1.50$mm。
求波长。
\exsolution
如存在未知微小中心膜厚 $h_0$，
同类暗环满足
$2h_0+r_k^2/R=k\lambda$。
两阶相减消去 $h_0$：
$(r_{15}^2-r_{10}^2)/R=5\lambda$。
因 $r_k=d_k/2$，
\[
\boxed{\lambda=
\frac{d_{15}^2-d_{10}^2}
{4R(15-10)}}
=\frac{(1.50^2-1.00^2)\times10^{-6}}
{4(0.125)(5)}
=\boxed{0.50\,\mathrm{\mu m}}.
\]
\excheck 单独用 $r_{10}^2=10\lambda R$ 与 $r_{15}^2=15\lambda R$ 会得到互相不一致的波长，说明应使用\emph{两环半径平方差}消去接触偏置。

\interferenceproblem{90}{201}
两片平凸薄玻璃透镜的球面相互接触，在反射光
$\lambda=0.60\,\mathrm{\mu m}$ 下，第五个亮牛顿环的直径为
$d=1.50$mm。
求两片薄透镜组成的光学系统的光焦度。
\exsolution
两个凸球面之间的空气隙近轴厚度
$h(r)=r^2(1/R_1+1/R_2)/2$，
令
$R_{\rm eff}^{-1}=R_1^{-1}+R_2^{-1}$，
第 $k$ 个亮环
$r_k^2=(k-\tfrac12)\lambda R_{\rm eff}$。
所以
\[
\frac1{R_{\rm eff}}
=\frac{(k-\tfrac12)\lambda}{r_k^2}
=\frac{4(k-\tfrac12)\lambda}{d^2}.
\]
每片平凸透镜的光焦度为
$(n-1)/R_i$，
接触总光焦度
$\Phi=(n-1)/R_{\rm eff}$，得
\[
\boxed{\Phi=
\frac{2(n-1)(2k-1)\lambda}{d^2}
\simeq2.4\,\mathrm D}
\]
（取玻璃 $n\simeq1.50$，$k=5$）。
\excheck 每片透镜一面为\emph{平面}，因此不能在透镜光焦度公式中凭空多乘球面系数2。

\interferenceproblem{91}{201}
两只对称薄玻璃透镜（一双凸、一双凹）贴合成光焦度
$\Phi=0.50$D 的系统，在空气隙中用
$\lambda=0.61\,\mathrm{\mu m}$ 的反射光观察牛顿环。
求：(a) 第十暗环的半径；
(b) 在空隙填水后该环半径如何变化。
\exsolution
两透镜曲面对应空气隙的等效曲率差
$1/R_{\rm gap}$；
因两透镜各有两面且对称，
其合光焦度满足
$\Phi=2(n_g-1)/R_{\rm gap}$。
因此
$R_{\rm gap}=2(n_g-1)/\Phi$。
暗环条件
$r_k^2=k\lambda R_{\rm gap}$，故
\[
\boxed{r_{10}=
\sqrt{\frac{2k\lambda(n_g-1)}{\Phi}}
\simeq3.5\,\mathrm{mm}}
\]
（$n_g\simeq1.50$）。
在水中膜光程差是原来的 $n_w$ 倍，
同一暗环
\[
\boxed{r'=\frac r{\sqrt{n_w}}
\simeq3.0\,\mathrm{mm}},\qquad n_w\simeq1.33.
\]
\excheck 两镜的净光焦度已知，避免单独求各面未知曲率；充水使相同阶次条纹的环半径变小。

\interferenceproblem{92}{201}
平凸透镜球面贴玻璃板，其空气隙换成二氧化碳介质（原书给
$n_1=1.50$ 为透镜、
$n_2=1.63$ 为填隙物、
$n_3=1.70$ 为玻璃板）。
球面半径 $R=100$cm，
反射光波长 $\lambda=0.50\,\mathrm{\mu m}$。
求第五暗环半径。
\exsolution
因 $n_1<n_2<n_3$，
在透镜—填隙物与填隙物—玻璃两界面反射均有 $\pi$ 相位反转，
相互反射波之间无额外 $\pi$ 相对翻转。
暗环条件为
$2n_2h=(k+\tfrac12)\lambda$，
$h=r^2/(2R)$，所以
\[
\boxed{r_k=
\sqrt{\frac{(2k+1)\lambda R}{2n_2}}}.
\]
对 $k=5$：
\[
\boxed{r_5\simeq
\sqrt{\frac{11(0.50\times10^{-6})(1)}
{2(1.63)}}
\simeq1.3\,\mathrm{mm}}.
\]
\excheck 本题与通常空气隙牛顿环的中心明暗相反，
必须用三个折射率确定两个反射相位。

\interferenceproblem{93}{202}
两光束干涉仪使用汞的橙色双线
$\lambda_1=576.97$nm、
$\lambda_2=579.03$nm。
求最早使两套干涉条纹彼此可见度最低的干涉级次。
\exsolution
两条谱线对应不同波数
$k_i=2\pi/\lambda_i$。
两套条纹的相位差相差 $\pi$ 时互相抵消，
因此
$2\pi\Delta(1/\lambda_1-1/\lambda_2)=\pi$。
以平均波长 $\bar\lambda$ 对应级次
$m=\Delta/\bar\lambda$，
有
\[
\boxed{m_{\min}\simeq
\frac{\bar\lambda}{2|\lambda_2-\lambda_1|}
=\frac{578\,\mathrm{nm}}{2(2.06\,\mathrm{nm})}
\simeq140}.
\]
\excheck 这里对应可见度包络的\emph{首次最小值}，不是一个单一波长的干涉强度零点。

\interferenceproblem{94}{202}
米歇尔逊干涉仪使用钠黄双线
$\lambda_1=589.0$nm、
$\lambda_2=589.6$nm。
解释为什么移动一镜时条纹会周期性消失，并求相邻两次最清晰干涉图样之间的镜面位移。
\exsolution
两波长的相对干涉相位
$\Delta\phi=2\pi\Delta(1/\lambda_1-1/\lambda_2)$。
因为各自干涉强度的周期稍不同，
当相对相位约为 $\pi$ 时相互抵消使条纹可见度最小；
相对相位增加 $2\pi$ 后两套条纹再度同相清晰。
相邻最清晰状态的光程差改变
$\delta\Delta=\lambda_1\lambda_2/|\lambda_2-\lambda_1|$。
镜移动 $\delta h$ 使往返光程差改变
$2\delta h$，所以
\[
\boxed{\delta h=
\frac{\lambda_1\lambda_2}{2|\lambda_2-\lambda_1|}
\simeq0.29\,\mathrm{mm}\approx0.30\,\mathrm{mm}}.
\]
\excheck 米歇尔逊反射为往返\emph{双程}，故镜面位移等于光程差增量的一半。

\interferenceproblem{95}{202}
厚度为 $d$ 的法布里—珀罗标准具受到单色光照射，在透镜焦平面出现同心干涉环（图5.18）。
求：(a) 各环位置随干涉级次的变化；
(b) 相邻极大角间距与干涉级次的关系。
\exsolution
标准具内光线与法线夹角为 $\theta$，
两次透射光相长的条件
\[
\boxed{2d\cos\theta=m\lambda},\qquad m\in\mathbb Z.
\]
令近轴透镜焦距为 $f$，
环半径 $r\simeq f\tan\theta\simeq f\theta$。
因 $\cos\theta$ 随 $\theta$ 增大而减小，
越向外的环对应\emph{更低}的 $m$。
相邻阶次差1，对上述式微分：
$2d|\sin\theta|\,\delta\theta\simeq\lambda$，
\[
\boxed{\delta\theta\simeq\frac{\lambda}{2d\sin\theta}}.
\]
\excheck 对较外侧较大 $\theta$，
同阶间隔在角度上\emph{减小}；在 $\theta=0$ 需用小角的平方关系不能直接除以 $\sin0$。

\interferenceproblem{96}{202}
法布里—珀罗标准具厚度 $d=2.5$cm，
观察 $\lambda=0.50\,\mathrm{\mu m}$ 的单色光。
求：(a) 最高干涉级次；
(b) 近此波长互不与相邻干涉级重叠的自由光谱范围 $\Delta\lambda$。
\exsolution
最大级次对应正入射，
$2d=m_{\max}\lambda$，
所以
\[
\boxed{m_{\max}=\frac{2d}{\lambda}
=\frac{0.050}{0.50\times10^{-6}}
=1.0\times10^5}.
\]
相邻谱线在不同干涉级次重叠时
$m\lambda=(m+1)(\lambda-\Delta\lambda)$，
对高阶 $m$，
$\Delta\lambda\simeq\lambda/m\simeq\lambda^2/(2d)$。
因此
\[
\boxed{\Delta\lambda\simeq
\frac{\lambda^2}{2d}
=\frac{(0.50\times10^{-6})^2}{0.050}
=5.0\,\mathrm{pm}}.
\]
\excheck 自由光谱范围与分辨本领的\emph{仪器精细度}不同；
前者由相邻级次重叠决定。
